
%We have implemented the \st{{\bf Plan Scoring} module} \matt{cost optimizer's {\bf Plan Scoring} module (discussed in \autoref{sec:planning})} through a number of techniques. The prototype contains a \st{database} \matt{catalog} of observations around model performance, cost, and quality collected from academic papers and service-published statistics. We use this \st{database} \matt{information}, along with \st{input data} statistics \matt{related to the size of the input relation}, to compute a very rough \matt{initial} score for each candidate plan. 

%The system also collects statistics about every plan during execution. If the plan is expected to run long enough, we will process a few inputs using preliminary "probe plans" that exist primarily to collect performance statistics. After the brief initial probe period, the system can then use these statistics to improve its estimates of the physical plan(s) it is executing \textbf{\textit{and}} other plans which it has not yet executed. For example, the distribution of source records' token lengths is useful information regardless of which model is chosen. On the other hand, some elements of cost and runtime cannot be estimated without switching physical plans.  



%This resembles a multi-armed bandit problem, where it may be beneficial to try multiple plans near the beginning of program execution in order to find the optimal plan, although we do not pursue this strategy in the current prototype.
%One particularly difficult challenge is estimating the quality of an operator's --- and a plan's --- generated outputs.  Currently we use both historical data and compare different models' outputs against those of a "champion" model that is known to have good quality (but which might have other unappealing characteristics).

%\system{} takes a user-written program and compiles it down to a (set of) physical plan(s) inside of its {\tt pz.Execute} class. In order to select one or more of the physical plans to execute, users must also provide a policy specifying their priorities regarding quality and cost, i.e., where they would like to operate on the pareto-frontier of plans in terms of the tradeoff between quality and cost. For example, a policy can minimize cost, maximize quality, or maximize the harmonic mean of cost and quality. Different policies may also allow the system to execute multiple plans with similar cost-quality tradeoffs in parallel to further maximize throughput.


%A crucial element of the \system{} system is its ability to re-optimize the physical plan(s) as it is executing a user's AI program. In brief, as \system{} executes a physical plan, it records the time and cost of each operation on a per-record basis. 


%This degree of caching presents two challenges. First, storing every intermediate result means that \system{} can quickly exceed the memory limits of even memory-optimized cloud compute instances. This implies that a significant portion of the cache must be spilled to disk. Second, for convert operations --- especially ones using prompt-level optimizations such as DSPy --- the continuous learning of an optimal (set of) prompt(s) means that the results of a convert operation may vary with each execution.

%Due to the orders of magnitude difference between the throughput of LLMs and every other component of the \system{} system, we are able to make unusual, yet reasonable design choices to mitigate these challenges. First, we recognize that reading from disk (or even remote storage) still provides orders of magnitude higher throughput (\textcolor{blue}{and cheaper cost}) than invoking an LLM to regenerate an output. Thus, spilling an entire history of intermediate results and reading them from disk can be a pareto-optimal design decision (on throughput and cost) depending on the user's workload. Second, for many workloads, small improvements in output generation quality may not be worth the cost of regenerating an entire \texttt{Dataset}. For example, if a user wants an estimate of an aggregate, it may be perfectly fine to re-use a cached computation even if a prompt optimizer learned an improved prompt. In each of these situations, the \system{} optimizer can decide whether to read from its cache or recompute an intermediate result based on the user workload, the optimizer's estimate of the cost vs. quality tradeoff, and the user's preferred policy (e.g., to minimize cost or maximize quality).
